\input{preamble}

\begin{document}

\title{Hodge Theory}
\author{Misha Verbitsky}
\date{IMPA, 2026}
\maketitle

\phantomsection
\label{section-phantom}

\tableofcontents

These notes follow the slides for Misha
Verbitsky's 2026 course at IMPA. The
course materials are available on the
\href{%
http://verbit.ru/IMPA/Kahler-2026/%
}{course page}.

\section{Lecture 1: complex manifolds}
\label{section-lecture-one-complex}

\noindent
This lecture was given on August 10,
2026. Erich K\"ahler described his work
as a golden thread leading into the
palpitating heart of mathematics.

\medskip\noindent
{\bf Complex structures.}

\begin{definition}
\label{definition-complex-operator}
Let $V$ be a real vector space. A
{\it complex structure operator} on $V$
is an automorphism
$$
I:V\longrightarrow V
$$
satisfying
$$
I^2=-\operatorname{Id}_V.
$$
\end{definition}

The action of $I$ extends
multiplicatively to every tensor space
built from $V$ and $V^*$. Thus, for
example,
$$
\begin{aligned}
I(&v_1\otimes\cdots\otimes v_k
\otimes w_1\otimes\cdots\otimes w_l)
\\
&=I(v_1)\otimes\cdots\otimes I(v_k)
\otimes I(w_1)\otimes\cdots\otimes
I(w_l).
\end{aligned}
$$

The following elementary facts will be
used repeatedly.

\begin{enumerate}
\item The eigenvalues of $I$ are
$\sqrt{-1}$ and $-\sqrt{-1}$.

\item There is an $I$-invariant positive
definite scalar product on $V$. Indeed,
starting with $g_0$, one may take
$$
g(x,y)=g_0(x,y)+g_0(Ix,Iy).
$$

\item The operator $I$ is orthogonal for
this scalar product, since
$$
g(Ix,Iy)=g(x,y).
$$

\item The operator $I$ is diagonalizable
over $\mathbb{C}$.

\item The two eigenvalues occur with the
same multiplicity.
\end{enumerate}

In a suitable complex basis of
$V\otimes_{\mathbb{R}}\mathbb{C}$, the
matrix of $I$ is
$$
\begin{bmatrix}
\sqrt{-1}\operatorname{Id}&0\\
0&-\sqrt{-1}\operatorname{Id}
\end{bmatrix}.
$$
In a suitable real basis it is block
diagonal, with every block equal to
$$
\begin{bmatrix}
0&-1\\
1&0
\end{bmatrix}.
$$

\begin{definition}
\label{definition-vector-hodge}
Let $(V,I)$ be a real vector space with a
complex structure. Its
{\it Hodge decomposition} is
$$
V\otimes_{\mathbb{R}}\mathbb{C}
=V^{1,0}\oplus V^{0,1},
$$
where $V^{1,0}$ and $V^{0,1}$ are the
$\sqrt{-1}$ and $-\sqrt{-1}$
eigenspaces of $I$, respectively.
\end{definition}

The dual space has the analogous Hodge
decomposition. Moreover, $V^{1,0}$
determines $I$: the operator acts by
$\sqrt{-1}$ on $V^{1,0}$ and by
$-\sqrt{-1}$ on its conjugate. We thus
obtain a bijection between complex
structures on $V$ and subspaces
$$
W\subset V\otimes_{\mathbb{R}}\mathbb{C}
$$
such that
$$
\dim_{\mathbb{C}}W
=\frac{1}{2}\dim_{\mathbb{R}}V,
\qquad
W\cap\overline{W}=0.
$$

\begin{definition}
\label{definition-hermitian-metric-vector}
An $I$-invariant positive definite scalar
product $g$ on $(V,I)$ is a
{\it Hermitian metric}. The triple
$(V,I,g)$ is a {\it Hermitian space}.
\end{definition}

If $g$ is Hermitian, the bilinear form
$$
\omega(x,y)=g(x,Iy)
$$
is skew-symmetric. Indeed,
$$
\begin{aligned}
\omega(x,y)
&=g(x,Iy)\\
&=g(Ix,I^2y)\\
&=-g(Ix,y)\\
&=-\omega(y,x).
\end{aligned}
$$

\begin{definition}
\label{definition-hermitian-form-vector}
The skew-symmetric form $\omega$ above is
the {\it Hermitian form} of $(V,I,g)$.
\end{definition}

Any two members of the triple
$(I,g,\omega)$ determine the third.

\medskip\noindent
{\bf Almost complex manifolds.}

\begin{definition}
\label{definition-almost-complex}
An {\it almost complex structure} on a
manifold $M$ is an endomorphism
$$
I:TM\longrightarrow TM
$$
such that
$$
I^2=-\operatorname{Id}_{TM}.
$$
The pair $(M,I)$ is an
{\it almost complex manifold}.
\end{definition}

On $\mathbb{C}^n$, write
$$
z_j=x_j+\sqrt{-1}y_j.
$$
The standard almost complex structure is
given by
$$
I\left(\frac{\partial}{\partial x_j}\right)
=\frac{\partial}{\partial y_j},
\qquad
I\left(\frac{\partial}{\partial y_j}\right)
=-\frac{\partial}{\partial x_j}.
$$

In real dimension two, an almost complex
structure is equivalent to an oriented
conformal structure.

\begin{exercise}
\label{exercise-almost-complex-conformal}
Prove this equivalence.
\end{exercise}

\begin{definition}
\label{definition-holomorphic-function-ac}
A function
$$
f:M\longrightarrow\mathbb{C}
$$
on an almost complex manifold is
{\it holomorphic} if
$$
df\in\Lambda^{1,0}(M).
$$
\end{definition}

Some almost complex manifolds have no
nonconstant holomorphic functions even
locally. A standard example is $S^6$ with
its $G_2$-invariant almost complex
structure.

\begin{theorem}
\label{theorem-holomorphic-equivalences}
Let $U\subset\mathbb{C}^n$ be open and let
$$
f:U\longrightarrow\mathbb{C}
$$
be differentiable. The following
conditions are equivalent.

\begin{enumerate}
\item The function $f$ is holomorphic.

\item At every point, the differential
$$
df:T_xU\longrightarrow\mathbb{C}
$$
is complex linear.

\item The restriction of $f$ to every
complex affine line is holomorphic as a
function of one complex variable.

\item Near every point, $f$ is represented
by a convergent complex Taylor series.
\end{enumerate}
\end{theorem}

\begin{proof}
The first two conditions are equivalent by
definition. A $(1,0)$-form restricts to a
$(1,0)$-form on every complex line. The
converse follows from the elementary fact
that a complex-valued linear form on
$(V,I)$ has type $(1,0)$ if and only if its
restriction to every $I$-invariant real
two-plane has type $(1,0)$. This proves the
equivalence with the third condition. A
convergent complex Taylor series is
holomorphic, while the converse follows
from the Cauchy formula.
\end{proof}

\begin{exercise}
\label{exercise-line-restrictions}
Prove the linear-algebraic fact used in the
proof of Theorem
\ref{theorem-holomorphic-equivalences}.
\end{exercise}

{\bf Sheaves and ringed spaces.}

\begin{definition}
\label{definition-presheaf-functions}
A {\it presheaf of functions} on a
topological space $M$ assigns a subring
$$
\mathcal{F}(U)\subset C(U)
$$
to every open set $U\subset M$, in such a
way that restriction from $U$ to any open
subset $U_1$ maps $\mathcal{F}(U)$ into
$\mathcal{F}(U_1)$.
\end{definition}

\begin{definition}
\label{definition-sheaf-functions}
A presheaf $\mathcal{F}$ is a
{\it sheaf of functions} if it also has the
following local-to-global property. Given
an open cover $U=\bigcup_iU_i$, a function
on $U$ belongs to $\mathcal{F}(U)$ whenever
its restriction to every $U_i$ belongs to
$\mathcal{F}(U_i)$.
\end{definition}

Continuous, smooth, and real analytic
functions form sheaves. Constant functions
and bounded functions form presheaves that
are not sheaves.

\begin{exercise}
\label{exercise-presheaves-not-sheaves}
Explain why constant functions and bounded
functions fail the sheaf condition.
\end{exercise}

\begin{definition}
\label{definition-ringed-space-course}
A {\it ringed space} $(M,\mathcal{F})$ is a
topological space with a sheaf of functions.
A {\it morphism of ringed spaces}
$$
\Psi:(M,\mathcal{F})
\longrightarrow(N,\mathcal{F}')
$$
is a continuous map such that, for every
open $U\subset N$ and every
$f\in\mathcal{F}'(U)$,
$$
\Psi^*f=f\circ\Psi
\in\mathcal{F}(\Psi^{-1}(U)).
$$
An {\it isomorphism of ringed spaces} is a
homeomorphism for which both directions are
morphisms.
\end{definition}

If $M$ is a $C^i$ manifold, the sheaf
$C^i$ makes $(M,C^i)$ a ringed space. Every
$C^i$ map induces a morphism between the
corresponding ringed spaces.

\begin{proposition}
\label{proposition-smooth-ringed-maps}
Morphisms between the ringed spaces
$(M,C^i)$ and $(N,C^i)$ are precisely the
$C^i$ maps from $M$ to $N$.
\end{proposition}

\begin{proof}
One implication follows because a
composition of $C^i$ functions is $C^i$.
Conversely, work in charts. A map between
open subsets of $\mathbb{R}^n$ is $C^i$ if
the pullback of every $C^i$ function is
$C^i$, since this condition applies in
particular to the coordinate functions.
\end{proof}

\begin{definition}
\label{definition-smooth-ringed-manifold}
A {\it smooth manifold of class $C^i$} is a
topological manifold with a sheaf of
functions which is locally isomorphic, as a
ringed space, to an open subset of
$(\mathbb{R}^n,C^i)$.
\end{definition}

\begin{definition}
\label{definition-chart-ringed}
A {\it chart}, or
{\it coordinate system}, on an open subset
of such a manifold is an isomorphism with
an open subset of $(\mathbb{R}^n,C^i)$.
\end{definition}

\begin{definition}
\label{definition-diffeomorphism-ringed}
A {\it diffeomorphism} is a homeomorphism
which induces an isomorphism of the
corresponding ringed spaces.
\end{definition}

These definitions agree with the usual
definitions in terms of atlases.

\begin{definition}
\label{definition-holomorphic-function-cn}
A {\it holomorphic function} on
$\mathbb{C}^n$ is a function whose
differential is complex linear, or
equivalently has type $(1,0)$.
\end{definition}

Holomorphic functions form a sheaf.

\begin{definition}
\label{definition-complex-manifold-ringed}
A {\it complex manifold} is a ringed space
locally isomorphic to an open ball in
$\mathbb{C}^n$ with its sheaf of
holomorphic functions.
\end{definition}

Equivalently, it is a manifold covered by
complex coordinate balls whose transition
functions are holomorphic.

\begin{definition}
\label{definition-holomorphic-ac-map}
A map
$$
\Psi:(M,I)\longrightarrow(N,J)
$$
of almost complex manifolds is
{\it holomorphic} if
$$
\Psi^*\Lambda^{1,0}(N)
\subset\Lambda^{1,0}(M).
$$
Equivalently, $d\Psi$ is complex linear.
\end{definition}

For the standard almost complex structures,
this is equivalent to holomorphicity of all
coordinate components of $\Psi$.

\begin{definition}
\label{definition-integrable-ac}
An almost complex structure $I$ on $M$ is
{\it integrable} if every point has a
neighborhood diffeomorphic to an open set
in $\mathbb{C}^n$ under a diffeomorphism
which identifies $I$ with the standard
almost complex structure.
\end{definition}

An integrable almost complex structure
defines a sheaf $\mathcal{O}_M$ of
holomorphic functions and makes
$(M,\mathcal{O}_M)$ a complex manifold.
Conversely, a complex manifold determines
an integrable almost complex structure. In
fact, its sheaf of holomorphic functions
determines the subbundle
$$
d\mathcal{O}_M
=\Lambda^{1,0}(M)
\subset\Lambda^1(M,\mathbb{C}),
$$
and this subbundle determines $I$.

\section{Assignment 1: complex structures}
\label{section-assignment-one-complex}

\noindent
This assignment was issued on August 10,
2026. Written solutions were requested for
the following Monday, with class discussion
on the Wednesday after that.

\begin{exercise}
\label{exercise-assignment-one-gl}
Prove that the group
$\operatorname{GL}(n,\mathbb{C})$ is
connected.
\end{exercise}

\begin{definition}
\label{definition-assignment-complex}
A {\it complex structure} on a real vector
space $V$ is an operator
$I\in\operatorname{End}(V)$ satisfying
$I^2=-\operatorname{Id}$.
\end{definition}

\begin{exercise}
\label{exercise-orthogonal-complex-four}
Let $V$ be a four-dimensional real vector
space with a Euclidean metric, and let $S$
be the space of orthogonal complex
structures on $V$. Prove that $S$ is a
disconnected union of two two-dimensional
spheres.
\end{exercise}

\begin{exercise}
\label{exercise-all-complex-components}
Let $V$ be an even-dimensional real vector
space, and let $S$ be the space of complex
structures on $V$. Prove that $S$ has two
connected components.
\end{exercise}

\begin{exercise}
\label{exercise-orthogonal-components}
Let $V$ be an even-dimensional real vector
space with a Euclidean metric $g$, and let
$S_g$ be the space of orthogonal complex
structures on $V$. Prove that $S_g$ has two
connected components.
\end{exercise}

\begin{exercise}
\label{exercise-two-forms-complex}
Let $V$ be a four-dimensional real vector
space, and let $U$ be the space of
skew-symmetric two-forms $\Omega$ on
$V\otimes_{\mathbb{R}}\mathbb{C}$ such that
$$
\Omega\wedge\Omega=0,
\qquad
\Omega\wedge\overline{\Omega}\ne0.
$$
Let $\mathbb{P}U$ be the projectivization
of $U$, namely the quotient of
$U\setminus\{0\}$ by the standard
$\mathbb{C}^*$-action. Construct a
$\operatorname{GL}(V)$-invariant bijection
between $\mathbb{P}U$ and the space of
complex structures on $V$.
\end{exercise}

\begin{proof}
Set $V_{\mathbb C}
=V\otimes_{\mathbb R}\mathbb C$.
We first recall that a nonzero two-form
$\Omega\in\Lambda^2V_{\mathbb C}^*$
is decomposable if and only if
$\Omega\wedge\Omega=0$.
Indeed, there exists a basis
$e^1,\ldots,e^4$ of $V_{\mathbb C}^*$
such that
\[
\Omega
=
\lambda_1e^1\wedge e^2
+
\lambda_2e^3\wedge e^4.
\]
Therefore
\[
\Omega\wedge\Omega
=
2\lambda_1\lambda_2
e^1\wedge e^2\wedge e^3\wedge e^4.
\]
Thus, if $\Omega\wedge\Omega=0$,
one of the $\lambda_i$ vanishes,
and hence $\Omega$ is decomposable.

Now let $\Omega\in U$. Write
$\Omega=\alpha\wedge\beta$ and define
\[
L_\Omega
=
\ker\Omega
=
\ker\alpha\cap\ker\beta.
\]
This is a two-dimensional complex
subspace of $V_{\mathbb C}$.
The condition
$\Omega\wedge\overline{\Omega}\ne0$
says that
$\alpha,\beta,\overline{\alpha},
\overline{\beta}$ are linearly
independent. Consequently,
\[
L_\Omega\cap\overline{L_\Omega}=0,
\qquad
V_{\mathbb C}
=
L_\Omega\oplus\overline{L_\Omega}.
\]

Define a complex-linear endomorphism of
$V_{\mathbb C}$ by
\[
J_\Omega|_{L_\Omega}
=
-i\,\operatorname{id},
\qquad
J_\Omega|_{\overline{L_\Omega}}
=
i\,\operatorname{id}.
\]
It commutes with complex conjugation,
so it is the complexification of a real
endomorphism $J_\Omega:V\to V$.
Moreover,
\[
J_\Omega^2=-\operatorname{id},
\]
so $J_\Omega$ is a complex structure
on $V$.

Since $L_{c\Omega}=L_\Omega$ for every
$c\in\mathbb C^*$, the complex structure
$J_\Omega$ depends only on the projective
class $[\Omega]$. Denote by
$\mathcal J(V)$ the set of complex
structures on $V$. We obtain a map
\[
\Phi:\mathbb PU\longrightarrow\mathcal J(V),
\qquad
[\Omega]\longmapsto J_\Omega.
\]

Conversely, let $J$ be a complex
structure on $V$. Its complexification
gives a decomposition
\[
V_{\mathbb C}
=
V_J^{1,0}\oplus V_J^{0,1},
\]
where $V_J^{1,0}$ and $V_J^{0,1}$ are
the $i$- and $-i$-eigenspaces,
respectively. Choose a nonzero form
\[
\Omega_J
\in
\Lambda^{2,0}V_{\mathbb C}^*.
\]
Equivalently, $\Omega_J$ is a nonzero
two-form with
$\ker\Omega_J=V_J^{0,1}$.
Since $V_J^{1,0}$ has complex dimension
two,
\[
\Omega_J\wedge\Omega_J=0.
\]
Furthermore, the decomposition of
$V_{\mathbb C}$ implies
\[
\Omega_J\wedge\overline{\Omega_J}\ne0.
\]
Thus $\Omega_J\in U$. Since
$\Lambda^{2,0}V_{\mathbb C}^*$ is
one-dimensional, the class $[\Omega_J]$
is independent of the choice of
$\Omega_J$.

These two constructions are inverse.
Finally, if $g\in\operatorname{GL}(V)$,
then
\[
\ker(g\cdot\Omega)=g(\ker\Omega),
\]
and hence
\[
J_{g\cdot\Omega}
=
gJ_\Omega g^{-1}.
\]
Therefore $\Phi$ is a
$\operatorname{GL}(V)$-equivariant
bijection.
\end{proof}


\begin{definition}
\label{definition-conformal-structure}
Two Riemannian metrics $g_1$ and $g_2$ on a
smooth manifold are {\it conformal}, or
{\it conformally equivalent}, if there is a
smooth positive function $\lambda$ such that
$$
g_1=\lambda g_2.
$$
A {\it conformal structure} is a metric up
to conformal equivalence.
\end{definition}

\begin{exercise}
\label{exercise-surface-conformal}
Let
$$
f:M\longrightarrow N
$$
be an orientation-preserving diffeomorphism
between almost Hermitian manifolds of real
dimension two. Prove that $f$ is
holomorphic if and only if it preserves the
conformal structures.
\end{exercise}

\begin{exercise}
\label{exercise-no-almost-complex}
Construct an oriented even-dimensional
manifold which admits no almost complex
structure.
\end{exercise}
\section{Lecture 2: real analytic
manifolds}
\label{section-lecture-two-real-analytic}

\noindent
This lecture was given on August 14, 2026.

\begin{definition}
\label{definition-real-analytic-function}
A {\it real analytic function} on an open
set $U\subset\mathbb{R}^n$ is a function
which has a convergent Taylor expansion near
every point. Thus, locally,
$$
\begin{aligned}
f(z_1+t_1,&\ldots,z_n+t_n)\\
&=\sum_{i_1,\ldots,i_n}
a_{i_1,\ldots,i_n}
t_1^{i_1}\cdots t_n^{i_n}.
\end{aligned}
$$
The real parameters $t_j$ are sufficiently
small for the expansion under consideration.
\end{definition}

Real analytic functions form a sheaf.

\begin{definition}
\label{definition-real-analytic-manifold}
A {\it real analytic manifold} is a ringed
space locally isomorphic to an open ball in
$\mathbb{R}^n$ equipped with its sheaf of
real analytic functions.
\end{definition}

Grauert's description of real analytic
manifolds in terms of complexifications is
developed below.

\medskip\noindent
{\bf Involutions and fixed points.}

\begin{definition}
\label{definition-involution}
An {\it involution} on $M$ is a map
$$
\iota:M\longrightarrow M
$$
satisfying
$$
\iota^2=\operatorname{Id}_M.
$$
\end{definition}

\begin{exercise}
\label{exercise-linear-involution}
Prove that every linear involution on a real
vector space is diagonalizable with
eigenvalues $1$ and $-1$.
\end{exercise}

\begin{theorem}
\label{theorem-fixed-set-smooth}
Let $M$ be a smooth manifold and let $\iota$
be a smooth involution. The fixed point set
of $\iota$ is a smooth submanifold.
\end{theorem}

\begin{proof}
Fix a point $m$ in the fixed point set. By
the inverse function theorem, $k$ functions
whose differentials are linearly independent
at $m$ define coordinates near $m$.

Suppose that $d\iota_m$ has $k$
eigenvalues equal to $1$ and $n-k$ equal to
$-1$. Choose coordinates $x_1,\ldots,x_n$
such that the first $k$ covectors at $m$ are
invariant and the remaining covectors are
anti-invariant. Set
$$
y_j=x_j+\iota^*x_j
$$
for $j\leq k$, and set
$$
y_j=x_j-\iota^*x_j
$$
for $j>k$. Up to harmless constant factors,
the $dy_j$ remain linearly independent at
$m$. They therefore give invariant local
coordinates, and the fixed point set is
locally cut out by
$$
y_{k+1}=\cdots=y_n=0.
$$
\end{proof}

\medskip\noindent
{\bf Real structures.}

\begin{definition}
\label{definition-real-structure-vector}
A {\it real structure} on a complex vector
space $V$ is an $\mathbb{R}$-linear
involution $\iota$ such that
$$
\iota(\lambda x)
=\overline{\lambda}\iota(x).
$$
\end{definition}

\begin{proposition}
\label{proposition-fixed-real-form}
Let $V^\iota$ be the fixed subspace of a
complex vector space with a real structure.
Then
$$
\dim_{\mathbb{R}}V^\iota
=\dim_{\mathbb{C}}V,
\qquad
V=V^\iota\otimes_{\mathbb{R}}\mathbb{C}.
$$
\end{proposition}

\begin{proof}
Suppose that $x_1,\ldots,x_n$ are linearly
independent in $V^\iota$ over $\mathbb{R}$.
Applying $\iota$ to a complex linear
relation among them and averaging shows that
the real parts of all coefficients vanish.
Applying the same argument after multiplying
the relation by $\sqrt{-1}$ shows that their
imaginary parts vanish. Hence the natural map
$$
V^\iota\otimes_{\mathbb{R}}\mathbb{C}
\longrightarrow V
$$
is injective.

It is also surjective. For $v\in V$, both
$$
\frac{v+\iota(v)}{2}
$$
and
$$
\frac{v-\iota(v)}{2\sqrt{-1}}
$$
belong to $V^\iota$, and together they
recover $v$ by complex linear combination.
\end{proof}

Let $\mathcal{C}_\iota$ be the category of
complex vector spaces with real structures,
where morphisms are complex linear maps
commuting with the involutions. Let
$\mathcal{C}_{\mathbb{R}}$ be the category
of real vector spaces.

\begin{theorem}
\label{theorem-real-structure-category}
The functor
$$
\begin{aligned}
\mathcal{C}_\iota
&\longrightarrow\mathcal{C}_{\mathbb{R}},\\
(V,\iota)&\longmapsto V^\iota
\end{aligned}
$$
is an equivalence of categories.
\end{theorem}

\begin{proof}
A quasi-inverse sends a real vector space
$W$ to
$$
W\otimes_{\mathbb{R}}\mathbb{C}
$$
with the involution
$$
w\otimes\lambda
\longmapsto
w\otimes\overline{\lambda}.
$$
The two composites are naturally isomorphic
to the identity by Proposition
\ref{proposition-fixed-real-form}.
\end{proof}

\begin{definition}
\label{definition-antiholomorphic-map}
A smooth map
$$
\Psi:(M,I)\longrightarrow(N,J)
$$
is {\it antiholomorphic} if
$$
d\Psi\circ I=-J\circ d\Psi.
$$
A complex-valued function is
{\it antiholomorphic} if its conjugate is
holomorphic.
\end{definition}

\begin{exercise}
\label{exercise-antiholomorphic-function}
Prove that an antiholomorphic function on
$M$ defines an antiholomorphic map from $M$
to $\mathbb{C}$.
\end{exercise}

\begin{exercise}
\label{exercise-antiholomorphic-pullback}
Prove that $\Psi$ is antiholomorphic if and
only if
$$
\Psi^*\Lambda^{0,1}(N)
\subset\Lambda^{1,0}(M).
$$
\end{exercise}

\begin{exercise}
\label{exercise-antiholomorphic-functions}
Let $\iota$ be a smooth self-map of a
complex manifold. Prove that $\iota$ is
antiholomorphic if and only if it pulls
every locally defined holomorphic function
back to an antiholomorphic one.
\end{exercise}

\begin{definition}
\label{definition-real-structure-manifold}
A {\it real structure} on a complex
manifold $M$ is an antiholomorphic
involution
$$
\iota:M\longrightarrow M.
$$
\end{definition}

\begin{proposition}
\label{proposition-real-local-coordinates}
Let $M$ be a complex manifold with a real
structure $\iota$. For every
$x\in M^\iota$, there is an
$\iota$-invariant coordinate neighborhood
with holomorphic coordinates
$z_1,\ldots,z_n$ such that
$$
\iota^*z_j=\overline{z_j}.
$$
\end{proposition}

\begin{proof}
Any basis of $\Lambda_x^{1,0}(M)$ occurs as
the differentials at $x$ of suitable
holomorphic coordinates: start from an
arbitrary chart and apply a complex linear
change of coordinates.

The differential $d\iota_x$ is a real
structure. A real basis of the cotangent
space to $M^\iota$ therefore determines a
complex basis of $T_x^*M$. Choose
holomorphic coordinates $u_1,\ldots,u_n$
with these prescribed differentials. The
functions
$$
z_j=u_j+\iota^*\overline{u_j}
$$
are holomorphic and satisfy the required
identity. After shrinking the neighborhood,
they form a coordinate system.
\end{proof}

\begin{definition}
\label{definition-compatible-real-atlas}
An atlas is {\it compatible with the real
structure} if every chart which meets the
fixed point set has the form given by
Proposition
\ref{proposition-real-local-coordinates}.
\end{definition}

\begin{proposition}
\label{proposition-fixed-real-analytic}
Let $M^\iota$ be the fixed point set of a
real structure on a complex manifold $M$.
A complex atlas compatible with the real
structure restricts to a real analytic
atlas on $M^\iota$.
\end{proposition}

\begin{proof}
The transition maps between compatible
charts commute with conjugation. They
therefore preserve the real locus and
restrict there to real-valued real analytic
maps.
\end{proof}

\begin{proposition}
\label{proposition-all-real-fixed-loci}
Every real analytic manifold arises as the
fixed point set of a real structure on a
complex manifold.
\end{proposition}

\begin{proof}
Take a locally finite real analytic atlas by
relatively compact balls $U_i$ with
transition maps $\Psi_{ij}$. For sufficiently
small $\varepsilon>0$, their convergent
Taylor series extend the transition maps to
complex neighborhoods
$$
\widetilde{U}_i
=U_i\times B_\varepsilon
\subset\mathbb{C}^n,
$$
where the second factor lies in the
imaginary directions. These complexified
transition maps glue the neighborhoods to a
complex manifold $M_{\mathbb{C}}$. Complex
conjugation is compatible with the gluing
and gives a real structure whose fixed point
set is the original manifold.
\end{proof}

\begin{definition}
\label{definition-complexification}
Let $M_{\mathbb{R}}$ be a real analytic
manifold. A {\it complexification} of
$M_{\mathbb{R}}$ is a complex manifold
$M_{\mathbb{C}}$ with a real structure whose
fixed point set is $M_{\mathbb{R}}$.
\end{definition}

There are several complexifications of the
circle. They include
$$
\mathbb{C}/\mathbb{Z}^2
$$
with $z\mapsto\overline z$, the space
$\mathbb{C}^*$ with
$$
z\longmapsto\overline z^{-1},
$$
and $\mathbb{C}P^1$ with
$$
[x:y]\longmapsto
[\overline{x}:\overline{y}].
$$
Their fixed loci are respectively
$\mathbb{R}/\mathbb{Z}$, the unit circle,
and $\mathbb{R}P^1$.

\begin{definition}
\label{definition-real-analytic-tensor}
A tensor on a real analytic manifold is
{\it real analytic} if, in real analytic
coordinates, it is a sum of coordinate
tensor monomials with real analytic
coefficients.
\end{definition}

\begin{definition}
\label{definition-same-germ}
Let $K\subset M$ and $K_1\subset M_1$ be
identified closed subsets of complex
manifolds. The manifolds $M$ and $M_1$
{\it have the same germ along $K$} if there
are neighborhoods $U$ and $U_1$ of the
identified set and a biholomorphism
$$
U\longrightarrow U_1
$$
which is the identity on $K$.
\end{definition}

\begin{definition}
\label{definition-germ-manifold}
The {\it germ of $M$ along $K$} is the
equivalence class of neighborhoods of $K$
under the relation in Definition
\ref{definition-same-germ}.
\end{definition}

Let $\mathcal{C}$ be the category of complex
manifolds with real structures and
holomorphic maps commuting with those
structures.

\begin{theorem}[Grauert]
\label{theorem-grauert-real-analytic}
The category of real analytic manifolds is
equivalent to the category of germs of
objects $(M,\iota)$ of $\mathcal{C}$ along
their fixed point sets $M^\iota$.
\end{theorem}

\begin{exercise}
\label{exercise-grauert-equivalence}
Prove Theorem
\ref{theorem-grauert-real-analytic}.
\end{exercise}

The theorem is associated with Hans Grauert,
who lived from 1930 to 2011.

\begin{definition}
\label{definition-antiholomorphic-tensor}
An {\it antiholomorphic tensor} on $(M,I)$
is a holomorphic tensor on $(M,-I)$.
\end{definition}

A function is holomorphic on $(M,I)$ if and
only if it is antiholomorphic on $(M,-I)$.
Complex conjugation identifies the sheaves
of holomorphic and antiholomorphic
functions. If $\iota$ is a real structure,
then $\iota^*$ sends holomorphic functions
and tensors to antiholomorphic ones.

\begin{theorem}
\label{theorem-extend-analytic-tensors}
Let $M_{\mathbb{C}}$ be a complexification
of $M_{\mathbb{R}}$. Every real analytic
tensor $\psi_{\mathbb{R}}$ extends to a
holomorphic tensor $\psi_{\mathbb{C}}$ on a
neighborhood of $M_{\mathbb{R}}$ such that
$$
\iota^*\psi_{\mathbb{C}}
=\overline{\psi_{\mathbb{C}}}.
$$
Conversely, every holomorphic tensor near
$M_{\mathbb{R}}$ satisfying this identity
restricts to a real analytic real tensor on
$M_{\mathbb{R}}$.
\end{theorem}

\begin{proof}
A holomorphic tensor restricts to a real
analytic tensor. The displayed identity
makes the restriction real. Conversely,
write a real analytic tensor in local
coordinates and replace the convergent real
Taylor series of its coefficients by their
complexifications. After shrinking the
neighborhoods, the resulting holomorphic
tensors agree on overlaps and glue.
\end{proof}

\medskip\noindent
\section{Lecture 2 bis: Frobenius theorem}
\label{section-lecture-two-bis-frobenius}

\noindent
This lecture was given on August 17, 2026.

\begin{definition}
\label{definition-distribution}
A {\it distribution}, or
{\it Pfaffian system}, on a manifold is a
subbundle
$$
B\subset TM.
$$
\end{definition}

Let
$$
\Pi:TM\longrightarrow TM/B
$$
be the quotient map. If $X$ and $Y$ are
sections of $B$, then
$$
[fX,Y]=f[X,Y]-Y(f)X.
$$
It follows that $\Pi([X,Y])$ is
$C^\infty(M)$-linear in both $X$ and $Y$.

\begin{definition}
\label{definition-frobenius-form}
The resulting bundle map
$$
\Lambda^2B\longrightarrow TM/B
$$
is the {\it Frobenius bracket}, or
{\it Frobenius form}, of $B$.
\end{definition}

\begin{definition}
\label{definition-involutive-distribution}
A distribution is {\it holonomic},
{\it involutive}, or {\it integrable} if
its Frobenius form vanishes.
\end{definition}

\medskip\noindent
{\bf Submersions and foliations.}

\begin{definition}
\label{definition-submersion-immersion}
A smooth map
$$
\pi:M\longrightarrow M'
$$
is a {\it submersion} if $d\pi$ is
surjective at every point, and an
{\it immersion} if $d\pi$ is injective at
every point.
\end{definition}

\begin{proposition}
\label{proposition-submersion-local-form}
Let $\pi:M\longrightarrow M'$ be a
submersion. Every point of $M$ has a
neighborhood $U$ isomorphic to a product
$V\times W$ in such a way that
$$
\pi|_U:V\times W\longrightarrow W
$$
is the second projection.
\end{proposition}

\begin{proof}
This is the local form supplied by the
inverse function theorem.
\end{proof}

\begin{definition}
\label{definition-vertical-tangent}
The {\it vertical tangent bundle} of a
submersion $\pi:M\longrightarrow M'$ is
$$
T_\pi M=\ker d\pi.
$$
\end{definition}

\begin{theorem}[Ehresmann]
\label{theorem-ehresmann-fibration}
A smooth proper submersion is a locally
trivial smooth fibration. In particular,
this holds for a smooth submersion between
compact manifolds.
\end{theorem}

\begin{theorem}[Frobenius]
\label{theorem-frobenius}
Let $B\subset TM$ be a subbundle. It is
involutive if and only if every point has a
neighborhood $U$ and a smooth submersion
$$
\pi:U\longrightarrow V
$$
such that
$$
B|_U=T_\pi U.
$$
\end{theorem}

One implication is immediate from the local
normal form for a submersion. The fibers of
$\pi$ are the local integral submanifolds of
$B$.

\begin{definition}
\label{definition-leaf-foliation}
A {\it leaf} of $B$ is a maximal connected
immersed submanifold tangent to $B$ at every
point. If $B$ is integrable, its collection
of leaves is the {\it foliation} associated
with $B$.
\end{definition}

Leaves need not be closed. To prove the
Frobenius theorem, it is enough to find an
integral submanifold through every point;
locally, the required submersion is then the
projection to the leaf space.

The problem was known as Pfaff's problem
before Cartan used the name Frobenius theorem
in 1922. Jacobi obtained an early weaker
result in 1837, Deahna supplied an important
technical argument in 1844, Clebsch treated
the generic case in 1866, and Frobenius gave
the general reformulation and solution in
1877. A historical account appears in
Thomas Hawkins's book on Frobenius in
context. The mathematicians highlighted in
the slides were Carl Gustav Jacobi
(1804--1851), Alfred Clebsch (1833--1872),
and Ferdinand Georg Frobenius (1849--1917).

\medskip\noindent
{\bf Ehresmann connections.}

\begin{definition}
\label{definition-ehresmann-connection}
Let $\pi:M\longrightarrow Z$ be a smooth
submersion. An {\it Ehresmann connection}
is a subbundle $T_{\rm hor}M\subset TM$
such that
$$
TM=T_{\rm hor}M\oplus T_\pi M.
$$
\end{definition}

Given a path
$$
\gamma:[0,a]\longrightarrow Z,
$$
parallel transport for an Ehresmann
connection is a smoothly varying family of
diffeomorphisms
$$
V_t:\pi^{-1}(\gamma(0))
\longrightarrow\pi^{-1}(\gamma(t))
$$
whose trajectories are tangent to
$T_{\rm hor}M$.

\begin{proposition}
\label{proposition-ehresmann-transport}
If $\pi$ is proper, parallel transport
exists along every compact path.
\end{proposition}

\begin{proof}
The horizontal lift is the solution of an
ordinary differential equation. Compactness
of the inverse image of the path prevents
the solution from escaping in finite time.
\end{proof}

This also gives the standard proof of
Theorem \ref{theorem-ehresmann-fibration}:
parallel transport along short radial paths
in a coordinate ball trivializes the
submersion over that ball.

\medskip\noindent
{\bf Proof of the Frobenius theorem.}

Suppose first that a Lie group $G$ acts on
$M$ and its infinitesimal vector fields span
a subbundle $B\subset TM$. The $G$-orbits
are integral submanifolds, so $B$ is
integrable. Also, if $X$ and $Y$ commute,
their local flows commute. This can be seen
in coordinates for which $X$ is a coordinate
vector field.

\begin{exercise}
\label{exercise-commuting-flows}
Prove that commuting vector fields have
commuting local flows.
\end{exercise}

Thus it is enough to construct locally a
basis of commuting vector fields for an
involutive distribution $B$. Fix a point
and choose a local submersion
$$
\sigma:M\longrightarrow M_1
$$
whose fibers are transverse to $B$. More
explicitly, in local coordinates take a
linear map to
$\mathbb{R}^{\operatorname{rank}B}$ whose
differential restricts to an isomorphism on
$B$ at the chosen point. After shrinking,
$$
d\sigma|_{B_x}:B_x
\longrightarrow T_{\sigma(x)}M_1
$$
is an isomorphism at every nearby point.

A vector field $X$ is projectable under
$\sigma$ if $d\sigma(X)$ is constant along
the fibers of $\sigma$. Projectable vector
fields are closed under the Lie bracket,
and
$$
d\sigma([X,Y])
=[d\sigma(X),d\sigma(Y)].
$$
Equivalently, if $\underline{X}$ is the
projected vector field, then
$$
\mathcal{L}_X(\sigma^*f)
=\sigma^*\mathcal{L}_{\underline{X}}f
$$
for every $f\in C^\infty(M_1)$.

Let $\zeta_1,\ldots,\zeta_k$ be coordinate
vector fields on $M_1$. There are unique
fields $\xi_1,\ldots,\xi_k$ in $B$ with
$$
d\sigma(\xi_i)=\zeta_i.
$$
They form a local basis of $B$. Since $B$
is involutive,
$$
[\xi_i,\xi_j]\in B.
$$
On the other hand,
$$
d\sigma([\xi_i,\xi_j])
=[\zeta_i,\zeta_j]=0.
$$
The restriction of $d\sigma$ to $B$ is
injective, so all these brackets vanish.
The commuting flows define a local action
whose orbits are the integral manifolds of
$B$, proving the theorem.

\section{Assignment 2: Frobenius form}
\label{section-assignment-two-frobenius}

\noindent
This assignment was issued on August 17,
2026, with the same submission and
discussion rules as Assignment 1.

\begin{exercise}
\label{exercise-dense-torus-flow}
Find a vector field on the two-dimensional
torus $T^2$ such that every orbit of its
flow is dense.
\end{exercise}

\begin{exercise}
\label{exercise-bracket-growth}
Construct a four-manifold $M$ and a rank
two distribution $B\subset TM$ such that
$[B,B]$ has rank three and $[[B,B],B]$ has
rank four.
\end{exercise}

\begin{definition}
\label{definition-contact-structure}
A {\it contact structure} on a
$(2n+1)$-dimensional manifold $M$ is a rank
$2n$ distribution $B\subset TM$ for which
$TM/B$ is a trivial line bundle and the
Frobenius form
$$
\Lambda^2B\longrightarrow TM/B
$$
is nondegenerate.
\end{definition}

\begin{exercise}
\label{exercise-rescale-contact-form}
Let $\theta$ be a one-form on a
$(2n+1)$-dimensional manifold such that
$$
\theta\wedge(d\theta)^n
$$
is nondegenerate. Let $f$ be a nowhere
vanishing smooth function and set
$\theta'=f\theta$. Prove that
$$
\theta'\wedge(d\theta')^n
$$
is nondegenerate.
\end{exercise}

\begin{exercise}
\label{exercise-kernel-contact}
Let $\theta$ be a one-form on a
$(2n+1)$-dimensional manifold such that
$\theta\wedge(d\theta)^n$ is nondegenerate.
Prove that
$$
\ker\theta\subset TM
$$
is a contact distribution.
\end{exercise}

\begin{exercise}
\label{exercise-contact-defining-form}
Let $M$ be an odd-dimensional manifold and
let $B\subset TM$ be a contact distribution.
Prove that there is a one-form $\theta$ such
that
$$
B=\ker\theta
$$
and $\theta\wedge(d\theta)^n$ is
nondegenerate.
\end{exercise}

\begin{exercise}
\label{exercise-contact-sphere}
Construct a contact structure on
$S^{2n+1}$ for every positive integer $n$.
\end{exercise}

\section{Lecture 3: the
Newlander--Nirenberg theorem}
\label{section-lecture-three-newlander}

\noindent
This lecture was given on August 17, 2026.
It begins with the real analytic manifolds,
real structures, compatible coordinates,
and complexifications constructed in
Lecture 2.

\begin{lemma}
\label{lemma-holomorphic-real-locus}
Let $X\subset\mathbb{C}^n$ be an open ball
invariant under conjugation, and let
$$
X_{\mathbb{R}}=X\cap\mathbb{R}^n.
$$
A holomorphic tensor on $X$ which vanishes
on $X_{\mathbb{R}}$ vanishes identically.
\end{lemma}

\begin{proof}
Every holomorphic coefficient which vanishes
on the real locus has all derivatives zero
there. Its Taylor series therefore vanishes,
and analytic continuation gives the result.
\end{proof}

\begin{definition}
\label{definition-real-analytic-ac}
An almost complex structure is
{\it real analytic} if its defining
endomorphism of $TM$ is a real analytic
tensor.
\end{definition}

\begin{lemma}[Corollary]
\label{lemma-complexified-ac-square}
Let $(M,I)$ be a real analytic almost
complex manifold. If $I_{\mathbb{C}}$ is
the holomorphic extension of $I$ to a
complexification, then
$$
I_{\mathbb{C}}^2
=-\operatorname{Id}.
$$
\end{lemma}

\begin{proof}
The holomorphic tensor
$$
I_{\mathbb{C}}^2+
\operatorname{Id}
$$
vanishes on the real locus, so Lemma
\ref{lemma-holomorphic-real-locus} applies.
\end{proof}

\medskip\noindent
{\bf The underlying real analytic
manifold.}

A holomorphic map between open subsets of
$\mathbb{C}^n$ is real analytic when viewed
as a map between open subsets of
$\mathbb{R}^{2n}$. Indeed, its real and
imaginary parts have convergent Taylor
series in the real variables.

\begin{definition}
\label{definition-underlying-real-analytic}
The {\it underlying real analytic manifold}
$M_{\mathbb{R}}$ of a complex manifold $M$
is the same manifold with its holomorphic
transition functions viewed as real analytic
maps.
\end{definition}

The real analytic functions on
$M_{\mathbb{R}}$ are the convergent power
series generated locally by holomorphic and
antiholomorphic functions. In holomorphic
coordinates, this is the same as taking
convergent power series in the real and
imaginary parts of the coordinates.

\begin{definition}
\label{definition-conjugate-manifold}
The {\it complex conjugate manifold}
$\overline{M}$ of $(M,I)$ is the same real
manifold equipped with $-I$. Its holomorphic
functions are the antiholomorphic functions
on $M$.
\end{definition}

\begin{proposition}
\label{proposition-product-complexification}
Let $M$ be a complex manifold. A
complexification of its underlying real
analytic manifold is
$$
M_{\mathbb{C}}=M\times\overline{M},
$$
with the real structure
$$
\tau(x,y)=(y,x).
$$
\end{proposition}

\begin{proof}
The fixed point set is the diagonal and is
naturally identified with
$M_{\mathbb{R}}$. Holomorphic functions from
the first factor and antiholomorphic
functions from the second restrict to the
corresponding functions on the diagonal.
Their convergent power series generate its
sheaf of real analytic functions.
\end{proof}

\medskip\noindent
{\bf Formal integrability.}

Recall that a complex manifold and an
integrable almost complex manifold are
equivalent notions, as explained after
Definition \ref{definition-integrable-ac}.

\begin{definition}
\label{definition-formally-integrable}
Let
$$
TM\otimes_{\mathbb{R}}\mathbb{C}
=T^{1,0}M\oplus T^{0,1}M
$$
be the Hodge decomposition of the
complexified tangent bundle. The almost
complex structure is
{\it formally integrable} if
$$
[T^{1,0}M,T^{1,0}M]
\subset T^{1,0}M.
$$
\end{definition}

\begin{definition}
\label{definition-nijenhuis-tensor}
The Frobenius form of $T^{1,0}M$, viewed as
a tensor measuring the component of
$[T^{1,0}M,T^{1,0}M]$ in $T^{0,1}M$, is the
{\it Nijenhuis tensor}.
\end{definition}

\begin{proposition}
\label{proposition-integrable-formal}
Every integrable almost complex structure is
formally integrable.
\end{proposition}

\begin{proof}
In holomorphic coordinates,
$T^{1,0}M$ is generated by the commuting
fields
$$
\frac{\partial}{\partial z_1},\ldots,
\frac{\partial}{\partial z_n}.
$$
Its Frobenius form therefore vanishes.
\end{proof}

\begin{theorem}[Newlander--Nirenberg]
\label{theorem-newlander-nirenberg}
An almost complex structure is integrable
if and only if it is formally integrable.
\end{theorem}

The name Nijenhuis tensor originates in a
1955 paper of Jan Schouten and Kentaro Yano
on the geometric meaning of its vanishing.
The slides also highlighted Albert Nijenhuis.
The real analytic version of the theorem was
proved in 1951 by Beno Eckmann and Alfred
Fr\"olicher. The latter surname is often
incorrectly written with an additional
letter.

The smooth theorem appeared in the sole
paper of August Newlander, written jointly
with Louis Nirenberg and published in 1957.
In a later interview Nirenberg recalled that
Andr\'e Weil and Shiing-Shen Chern drew his
attention to the problem. The task was to
recognize Cauchy--Riemann operators in
arbitrary coordinates. The first approach
removed error terms one at a time and worked
in complex dimension two, but failed in
higher dimensions. The eventual proof used
a fully nonlinear change of variables to
remove all error terms simultaneously.

\medskip\noindent
{\bf The real analytic case.}

Let $(M,I)$ be real analytic and formally
integrable. Extend $I$ holomorphically to a
complexification $M_{\mathbb{C}}$. After
shrinking, Lemma
\ref{lemma-complexified-ac-square} gives a
holomorphic decomposition
$$
TM_{\mathbb{C}}
=T^{1,0}M_{\mathbb{C}}
\oplus T^{0,1}M_{\mathbb{C}}.
$$
Formal integrability on the real locus and
Lemma \ref{lemma-holomorphic-real-locus}
show that both distributions are involutive
near that locus. The Frobenius theorem gives
their $(1,0)$ and $(0,1)$ foliations.

When $I$ is already integrable, one may use
the complexification
$$
M_{\mathbb{C}}=M\times\overline{M}.
$$
If $\pi$ and $\overline{\pi}$ are the two
projections, the fibers of
$\overline{\pi}$ are the $(1,0)$ leaves and
the fibers of $\pi$ are the $(0,1)$ leaves.
Along the diagonal, this splitting is the
usual Hodge splitting. It follows in
particular that an integrable almost complex
structure is real analytic with respect to
its underlying real analytic atlas.

Holomorphic functions on
$M_{\mathbb{C}}$ which are constant along
the $(0,1)$ leaves restrict to holomorphic
functions on $(M,I)$, because their
differentials vanish in every $(0,1)$
direction.

\begin{theorem}
\label{theorem-real-analytic-newlander}
Every real analytic, formally integrable
almost complex structure is integrable.
\end{theorem}

\begin{proof}
The two foliations above are transverse.
Consequently, the complexification is
locally a product of their leaf spaces,
say
$$
M_{\mathbb{C}}=M'\times M''.
$$
Functions on $M'$ lift to functions on the
product which are constant along the
$(0,1)$ leaves. Choose
$$
n=\frac{1}{2}\dim_{\mathbb{R}}M
$$
such functions $f_1,\ldots,f_n$ whose
differentials are linearly independent at a
given point. Their restrictions to $M$ are
holomorphic. The inverse function theorem
shows that they form local complex
coordinates, and their transition functions
are holomorphic by construction.
\end{proof}

\section{Assignment 3: almost complex
structures and integrability}
\label{section-assignment-three-integrable}

\noindent
This assignment was issued on August 24,
2026, with the same submission and
discussion rules as the earlier assignments.

\begin{definition}
\label{definition-hodge-component-operator}
Let $(M,I)$ be an almost complex manifold
and let
$$
A:\Lambda^*M\longrightarrow\Lambda^*M
$$
be linear. The {\it Hodge components} of
$A$ are the operators $A^{p,q}$ such that
$$
A=\sum_{p,q}A^{p,q}
$$
and
$$
A^{p,q}(\Lambda^{i,j}(M))
\subset\Lambda^{i+p,j+q}(M).
$$
\end{definition}

\begin{exercise}
\label{exercise-four-hodge-components}
Prove that the de Rham differential on an
almost complex manifold has at most four
nonzero Hodge components:
$$
d=d^{2,-1}+d^{1,0}+d^{0,1}+d^{-1,2}.
$$
\end{exercise}

\begin{exercise}
\label{exercise-cartan-formulas}
Let $\theta$ be a one-form on a smooth
manifold $M$, and let $X,Y$ be vector
fields.

\begin{enumerate}
\item Let
$$
\iota_X:\Lambda^i(M)
\longrightarrow\Lambda^{i-1}(M)
$$
be contraction with $X$. Prove Cartan's
formula
$$
\mathcal{L}_X=\iota_Xd+d\iota_X.
$$

\item Prove that
$$
\begin{aligned}
X(\theta(Y))
&=\theta([X,Y])\\
&\quad+Y(\theta(X))\\
&\quad+d\theta(X,Y).
\end{aligned}
$$

\item Deduce the equivalent identity
$$
\begin{aligned}
Y(\theta(X))
&=-\theta([X,Y])\\
&\quad+X(\theta(Y))\\
&\quad-d\theta(X,Y).
\end{aligned}
$$

\item Prove the formula
$$
\begin{aligned}
d\theta(X,Y)
&=X(\theta(Y))-Y(\theta(X))\\
&\quad-\theta([X,Y]).
\end{aligned}
$$
\end{enumerate}
\end{exercise}

For the first part, use the general
properties of derivations. Deduce the later
identities successively from Cartan's
formula.

\begin{exercise}
\label{exercise-integrability-components}
Let $(M,I)$ be an almost complex manifold
and write
$$
d=d^{2,-1}+d^{1,0}+d^{0,1}+d^{-1,2}.
$$

\begin{enumerate}
\item Prove that $d^{2,-1}$ and
$d^{-1,2}$ are $C^\infty(M)$-linear.

\item Suppose that
$$
d^{2,-1}=d^{-1,2}=0.
$$
Prove that $I$ is formally integrable.

\item Suppose that
$$
(d^{0,1})^2=0.
$$
Prove that $I$ is formally integrable.
\end{enumerate}
\end{exercise}

Use the formula in Exercise
\ref{exercise-cartan-formulas} for the
second part and $d^2=0$ for the third.

\begin{exercise}
\label{exercise-nijenhuis-hodge-component}
Let $(M,I)$ be an almost complex manifold.

\begin{enumerate}
\item Prove that
$$
d^{-1,2}:\Lambda^{1,0}(M)
\longrightarrow\Lambda^{0,2}(M)
$$
is dual to the conjugate Nijenhuis tensor
$$
\overline{N}:\Lambda^2(T^{0,1}M)
\longrightarrow T^{1,0}M.
$$

\item Prove that
$$
d^{-1,2}(d\varphi)=0
$$
for every holomorphic function $\varphi$.
\end{enumerate}
\end{exercise}

\begin{exercise}
\label{exercise-surjective-nijenhuis}
Let $(M,I)$ be an almost complex manifold.
Suppose that the Nijenhuis tensor
$$
N:\Lambda^2(T^{1,0}M)
\longrightarrow T^{0,1}M
$$
is surjective. Prove that $(M,I)$ has no
nonconstant local holomorphic functions.
\end{exercise}

Use Exercise
\ref{exercise-nijenhuis-hodge-component}.

\begin{definition}
\label{definition-left-invariant-complex}
Let $G$ be a real Lie group with Lie algebra
$\mathfrak{g}=T_eG$. The {\it left action}
of $G$ on itself is
$$
L_g(x)=gx.
$$
An almost complex structure on $G$ is
{\it left invariant} if it is preserved by
$L_g$ for every $g\in G$.
\end{definition}

\begin{exercise}
\label{exercise-left-invariant-integrable}
Let $G$ be a real Lie group, let
$\mathfrak{g}=T_eG$, and set
$$
\mathfrak{g}_{\mathbb{C}}
=\mathfrak{g}\otimes_{\mathbb{R}}
\mathbb{C}.
$$

\begin{enumerate}
\item Prove that every complex structure
operator on $\mathfrak{g}$ extends uniquely
to a left-invariant almost complex structure
on $G$.

\item Let $I$ be a complex structure on
$\mathfrak{g}$ and write
$$
\mathfrak{g}_{\mathbb{C}}
=\mathfrak{g}^{1,0}
\oplus\mathfrak{g}^{0,1}.
$$
Prove that the induced left-invariant
structure is formally integrable if and
only if
$$
\mathfrak{g}^{1,0}
\subset\mathfrak{g}_{\mathbb{C}}
$$
is a Lie subalgebra.
\end{enumerate}
\end{exercise}

\begin{exercise}[Optional]
\label{exercise-left-invariant-no-functions}
Construct a left-invariant almost complex
structure $I$ on a Lie group $G$ such that
$(G,I)$ has no nonconstant local
holomorphic functions.
\end{exercise}

Use Exercise
\ref{exercise-surjective-nijenhuis}.

\begin{exercise}
\label{exercise-holomorphic-generators}
Let $(M,I)$ be an almost complex manifold.
Suppose that $\Lambda^{1,0}(M)$ is generated
by differentials of holomorphic functions.
Without using Theorem
\ref{theorem-newlander-nirenberg}, prove
that $I$ is integrable.
\end{exercise}

\section{Lecture 4: K\"ahler manifolds}
\label{section-lecture-four-kahler}

\noindent
Note: in this lecture we proved
the existence on Fubini-Study metric
on $\mathbb{C}P^n$ (and apparently also
in any symmetric space)
using $U(n+1)$-invariance.
Namely, a theorem states that there
is a bijection between $G$-invariant
metrics on $M$ and
$\operatorname{St}_x(G)$-invariant
metrics on $T_xM$.

For symmetric spaces, there are no
odd invariant tensors,
but $d \omega$ would be such.
And thus $\omega$ is closed.

We also proved that a complex
submanifold of a Kähler manifold is also
Kähler.

\medskip\noindent
{\bf }This lecture was given on August 21,
2026.

\medskip\noindent
{\bf Homogeneous spaces.}

\begin{definition}
\label{definition-lie-group-action}
A {\it Lie group} is a smooth manifold with
a group structure for which multiplication
and inversion are smooth. A
{\it smooth action} of a Lie group $G$ on a
manifold $M$ is a smooth map
$$
G\times M\longrightarrow M
$$
which is a group action.
\end{definition}

\begin{definition}
\label{definition-homogeneous-space}
A $G$-manifold $M$ is a
{\it homogeneous space} if the action is
transitive. For $x\in M$, the
{\it stabilizer}, or
{\it isotropy subgroup}, is
$$
\operatorname{Stab}_x(G)
=\{g\in G\mid g(x)=x\}.
$$
\end{definition}

\begin{proposition}
\label{proposition-homogeneous-quotient}
If $G$ acts transitively on $M$, then
$$
M\cong G/H,
$$
where $H=\operatorname{Stab}_x(G)$ for any
chosen point $x\in M$.
\end{proposition}

\begin{proof}
The surjective orbit map
$$
G\longrightarrow M,
\qquad
g\longmapsto g(x),
$$
has fibers equal to the left cosets of $H$.
\end{proof}

If $g(x)=y$, the stabilizers of $x$ and $y$
are conjugate.

\begin{definition}
\label{definition-isotropy-representation}
The {\it isotropy representation} at
$x\in G/H$ is the natural action of
$\operatorname{Stab}_x(G)$ on $T_xM$.
\end{definition}

\begin{definition}
\label{definition-invariant-tensor}
A tensor on a homogeneous space $G/H$ is
{\it invariant} if it is preserved by the
diffeomorphisms induced by all elements of
$G$.
\end{definition}

An isotropy-invariant tensor $\Phi_x$ on
$T_xM$ can be transported to any
$y=g(x)$ by the action of $g$. This is
independent of the choice of $g$, because
two choices differ by an element of the
stabilizer.

\begin{theorem}
\label{theorem-invariant-isotropy-tensors}
For a homogeneous space $M=G/H$, invariant
tensors on $M$ are in bijection with
isotropy-invariant tensors on $T_xM$.
\end{theorem}

\medskip\noindent
{\bf Hermitian and K\"ahler manifolds.}

\begin{definition}
\label{definition-hermitian-metric-manifold}
A Riemannian metric $g$ on an almost complex
manifold is {\it Hermitian} if
$$
g(IX,IY)=g(X,Y).
$$
\end{definition}

Every Riemannian metric $g$ produces a
Hermitian metric
$$
h(X,Y)=g(X,Y)+g(IX,IY).
$$

\begin{definition}
\label{definition-hermitian-form-manifold}
The {\it Hermitian form} of an almost
Hermitian manifold $(M,I,g)$ is
$$
\omega(X,Y)=g(X,IY).
$$
\end{definition}

The Hermitian form is skew-symmetric and
has Hodge type $(1,1)$. As in the vector
space case, any two members of
$(I,g,\omega)$ determine the third.

\begin{definition}
\label{definition-kahler-manifold}
A complex Hermitian manifold $(M,I,g)$ is
{\it K\"ahler} if its Hermitian form is
closed. In this case $\omega$ is the
{\it K\"ahler form}, and
$$
[\omega]\in H^2(M,\mathbb{R})
$$
is its {\it K\"ahler class}.
\end{definition}

The slides recalled Erich K\"ahler
(1906--2000) and Andr\'e Weil
(1906--1998), including a childhood
photograph of Andr\'e with Simone Weil.

\medskip\noindent
{\bf Invariant metrics.}

\begin{theorem}
\label{theorem-transitive-sphere-form}
Let $G$ act linearly on $V$ and transitively
on the unit sphere of a positive definite
quadratic form $g$. Every $G$-invariant
symmetric bilinear form on $V$ is a scalar
multiple of $g$.
\end{theorem}

\begin{proof}
The form $g$ is invariant because $G$
preserves its unit sphere. If $g'$ is
another invariant quadratic form, then
$$
x\longmapsto\frac{g'(x,x)}{g(x,x)}
$$
is constant on the unit sphere by
transitivity and invariant under nonzero
real scaling. It is therefore constant on
$V\setminus\{0\}$.
\end{proof}

\begin{exercise}
\label{exercise-transitive-irreducible}
Let $V$ be a representation of $G$. Suppose
that $G$ acts transitively on a sphere in
$V$. Prove that $V$ is irreducible.
\end{exercise}

\begin{exercise}
\label{exercise-real-schur-form}
Prove the following real form of Schur's
lemma. If $V$ is an irreducible real
representation of $G$ and $g$ is an
invariant positive definite symmetric
bilinear form, then every invariant
symmetric bilinear form is proportional to
$g$.
\end{exercise}

\medskip\noindent
{\bf The Fubini--Study metric.}

The unitary group $U(n+1)$ acts transitively
on $\mathbb{C}P^n$, and the stabilizer of a
point is $U(n)\times U(1)$.

\begin{theorem}
\label{theorem-fubini-study}
There is a $U(n+1)$-invariant Riemannian
metric on $\mathbb{C}P^n$. It is unique up
to a positive scalar and is K\"ahler.
\end{theorem}

\begin{proof}
At a point, the tangent space is the standard
representation of $U(n)$. Its standard
positive form is isotropy invariant and
therefore extends to an invariant metric by
Theorem
\ref{theorem-invariant-isotropy-tensors}.
The isotropy group acts transitively on the
unit sphere, so uniqueness follows from
Theorem \ref{theorem-transitive-sphere-form}.

Let $\omega$ be the Hermitian form of this
metric. The three-form $d\omega$ is
isotropy invariant. Since the isotropy group
contains $-\operatorname{Id}$, every
invariant tensor of odd covariant degree
vanishes. Thus $d\omega=0$.
\end{proof}

\begin{definition}
\label{definition-fubini-study}
The invariant metric in Theorem
\ref{theorem-fubini-study} is the
{\it Fubini--Study metric}; its Hermitian
form is the {\it Fubini--Study form}.
\end{definition}

\medskip\noindent
{\bf Projective manifolds.}

\begin{definition}
\label{definition-complex-submanifold}
A smooth submanifold $X\subset M$ of a
complex manifold is a
{\it complex submanifold} if
$$
I(TX)\subset TX.
$$
The inclusion is then a
{\it complex embedding}.
\end{definition}

\begin{definition}
\label{definition-projective-manifold}
A complex manifold is {\it projective} if
it admits a complex embedding into some
$\mathbb{C}P^n$.
\end{definition}

A complex submanifold of a K\"ahler manifold
is K\"ahler: restrict both the Hermitian
metric and its closed Hermitian form. Hence
every projective manifold is K\"ahler.

\begin{definition}
\label{definition-projective-algebraic}
A subvariety of $\mathbb{C}P^n$ is
{\it complex algebraic} if it is the common
zero locus of homogeneous polynomials.
\end{definition}

\begin{theorem}[Chow]
\label{theorem-chow-projective}
Every complex submanifold of
$\mathbb{C}P^n$ is complex algebraic.
\end{theorem}

\begin{definition}
\label{definition-integral-kahler-class}
A K\"ahler manifold has an
{\it integral K\"ahler class} if its K\"ahler
class lies in the image of
$$
H^2(M,\mathbb{Z})
\longrightarrow H^2(M,\mathbb{R}).
$$
\end{definition}

Since
$$
H^2(\mathbb{C}P^n,\mathbb{R})
\cong\mathbb{R},
$$
the Fubini--Study form can be scaled to
have integral cohomology class. Its
restriction shows that every projective
manifold has a K\"ahler metric with integral
K\"ahler class.

\begin{theorem}[Kodaira embedding]
\label{theorem-kodaira-embedding}
Every compact K\"ahler manifold with an
integral K\"ahler class is projective.
\end{theorem}

The slides conclude with the hierarchy of
classes of almost complex manifolds. Every
projective manifold is both K\"ahler and
Moishezon. Every K\"ahler manifold is both
complex and symplectic. Every Moishezon
manifold is complex. Finally, complex and
symplectic manifolds both carry almost
complex structures. The slides refer to
Moishezon manifolds as algebraic spaces.

\section{Lecture 5: the Levi--Civita
connection}
\label{section-lecture-five-levi-civita}

\noindent
Maybe the highlight of this lecture
(to which I arrived late)
was the proof of the existence
and uniqueness of the LC connection
via a non-conventional proof:
existence of metric connections is easy,
and each is determined by its torsion
via an isomorphism proved
combinatorially, which is the heart
of the proof. Taking preimage of zero
under such isomorphism we obtain
a metric torsion-free connection.

\medskip\noindent
This lecture was given on August 24, 2026.

\medskip\noindent
{\bf Connections and parallel transport.}

\begin{definition}
\label{definition-vector-bundle-connection}
A {\it connection} on a vector bundle $B$
over $M$ is an operator
$$
\nabla:\Gamma(B)
\longrightarrow
\Gamma(B\otimes\Lambda^1M)
$$
satisfying
$$
\nabla(fb)=b\otimes df+f\nabla b.
$$
For a vector field $X$, contraction with
$X$ gives the covariant derivative
$\nabla_Xb$.
\end{definition}

A connection on $B$ induces a connection on
$B^*$, characterized by
$$
\begin{aligned}
d\left\langle{}b,\beta\right\rangle{}
&=\left\langle{}\nabla b,
\beta\right\rangle{}\\
&\quad+
\left\langle{}b,
\nabla\beta\right\rangle{}.
\end{aligned}
$$
The same symbol $\nabla$ is conventionally
used for both. The Leibniz rule
$$
\begin{aligned}
\nabla(b_1\otimes b_2)
&=(\nabla b_1)\otimes b_2\\
&\quad+b_1\otimes\nabla b_2
\end{aligned}
$$
then defines connections on all tensor
bundles built from $B$ and $B^*$.

\begin{theorem}
\label{theorem-parallel-section-line}
Let $B$ be a vector bundle with connection
over the real line. For each $x$ and each
$b_x\in B_x$, there is a unique parallel
section $b$ satisfying
$$
\nabla b=0,
\qquad
b|_x=b_x.
$$
\end{theorem}

\begin{proof}
In a trivialization this is the existence
and uniqueness theorem for the ordinary
differential equation
$$
\frac{db}{dt}+A(t)b=0.
$$
\end{proof}

\begin{definition}
\label{definition-parallel-transport}
Let
$$
\gamma:[0,1]\longrightarrow M
$$
be a smooth path from $x$ to $y$. The
{\it parallel transport} of $b_x\in B_x$
along $\gamma$ is the vector $b_y\in B_y$
obtained by solving
$$
\nabla_{\dot\gamma}b=0,
\qquad
b(0)=b_x,
$$
and taking $b_y=b(1)$.
\end{definition}

\begin{definition}
\label{definition-holonomy-group}
The {\it holonomy group} of a connection at
$x$ is the group of automorphisms of $B_x$
obtained by parallel transport around all
loops based at $x$.
\end{definition}

\medskip\noindent
{\bf Lie algebras and tensors.}

If $V$ is a representation of a Lie algebra
$\mathfrak{g}$, then $V^*$ becomes a
representation by
$$
\left\langle{}a(x),\lambda
\right\rangle{}
=-
\left\langle{}x,a(\lambda)
\right\rangle{}.
$$
Tensor products become representations by
the Leibniz rule
$$
a(x\otimes y)
=a(x)\otimes y+x\otimes a(y).
$$
Thus $\mathfrak{g}$ acts on all tensor
powers of $V$ and $V^*$.

\begin{definition}
\label{definition-lie-preserves-tensor}
The Lie algebra $\mathfrak{g}$
{\it preserves a tensor} $\Phi$ if
$$
a(\Phi)=0
$$
for every $a\in\mathfrak{g}$.
\end{definition}

\begin{definition}
\label{definition-orthogonal-lie-algebra}
Let $h$ be a nondegenerate symmetric
bilinear form on $V$. The
{\it orthogonal Lie algebra}
$\mathfrak{so}(V,h)$ consists of the
endomorphisms $a$ satisfying
$$
h(a(x),y)=-h(x,a(y)).
$$
\end{definition}

In an $h$-orthonormal basis these are the
skew-symmetric matrices.

\begin{proposition}
\label{proposition-two-forms-so}
Use $h$ to identify $V$ with $V^*$ and let
$$
\tau:V^*\otimes V^*
\longrightarrow V^*\otimes V
$$
be the resulting isomorphism. Then
$$
\tau(\Lambda^2V^*)
=\mathfrak{so}(V,h).
$$
\end{proposition}

\begin{proof}
The bilinear form corresponding to
$a\in\operatorname{End}(V)$ is
$$
(x,y)\longmapsto h(a(x),y).
$$
It is skew-symmetric precisely when $a$
satisfies the defining identity for
$\mathfrak{so}(V,h)$.
\end{proof}

\begin{definition}
\label{definition-parallel-tensor}
A tensor $\Psi$ built from a vector bundle
$B$ and its dual is {\it parallel}, or
{\it preserved by $\nabla$}, if
$$
\nabla\Psi=0.
$$
\end{definition}

This condition holds if and only if every
parallel transport preserves $\Psi$.

\begin{lemma}[Corollary]
\label{lemma-parallel-holonomy}
A tensor is parallel if and only if it is
invariant under the holonomy group.
\end{lemma}

\begin{definition}
\label{definition-orthogonal-connection}
Let $h$ be a positive definite form on $B$.
A connection is {\it orthogonal} if
$$
\nabla h=0.
$$
\end{definition}

\begin{definition}
\label{definition-unitary-connection}
Let $(B,I,h)$ be a Hermitian vector bundle.
A connection is {\it unitary}, or
{\it Hermitian}, if
$$
\nabla I=0,
\qquad
\nabla h=0.
$$
\end{definition}

\medskip\noindent
{\bf Torsors and affine spaces.}

\begin{definition}
\label{definition-torsor}
A {\it torsor} over a group $G$ is a set
$X$ with a free and transitive $G$-action.
\end{definition}

\begin{definition}
\label{definition-torsor-morphism}
A {\it morphism of torsors} from
$(X,G,\rho)$ to $(X',G',\rho')$ is a pair
of maps
$$
\Psi_X:X\longrightarrow X',
\qquad
\Psi_G:G\longrightarrow G',
$$
where $\Psi_G$ is a homomorphism and
$$
\Psi_X(\rho(g,x))
=\rho'(\Psi_G(g),\Psi_X(x)).
$$
\end{definition}

These objects and morphisms form the
category of torsors.

\begin{definition}
\label{definition-affine-space}
An {\it affine space} is a torsor over a
vector space $V$. The vector space $V$ is
its {\it linearization}.
\end{definition}

If $\nabla$ and $\nabla_1$ are connections
on $B$, their difference is an
$\operatorname{End}(B)$-valued one-form.
Conversely, adding any
$$
A\in\Lambda^1M\otimes
\operatorname{End}(B)
$$
to a connection produces another
connection. Hence the space of connections
on $B$ is an affine space over
$$
\Lambda^1M\otimes\operatorname{End}(B).
$$

\begin{proposition}
\label{proposition-affine-orthogonal}
Let $B$ carry a scalar product. The space of
orthogonal connections on $B$ is an affine
space over
$$
\Lambda^1M\otimes\mathfrak{so}(B).
$$
\end{proposition}

\begin{proof}
The difference $A$ of two connections acts
on a bilinear form $h$ by
$$
\begin{aligned}
(A(h))(x,y)
&=-h(A(x),y)\\
&\quad-h(x,A(y)).
\end{aligned}
$$
Thus two orthogonal connections differ by a
one-form with values in
$\mathfrak{so}(B)$, and every such change
preserves orthogonality.
\end{proof}

More generally, let $B$ carry a tensor
$\Phi$ and let
$$
\mathfrak{g}\subset\operatorname{End}(B)
$$
be the Lie algebra preserving $\Phi$. The
space of connections preserving $\Phi$ is
an affine space over
$$
\Lambda^1M\otimes\mathfrak{g}.
$$

\medskip\noindent
{\bf The de Rham differential.}

\begin{definition}
\label{definition-de-rham-algebra-course}
The {\it de Rham algebra} of $M$ is the
graded algebra
$$
\Lambda^*M
=\Gamma(\Lambda^*T^*M)
$$
with the wedge product. Its degree $i$
elements are the
{\it differential $i$-forms}.
\end{definition}

In degree zero one has
$$
\Lambda^0M=C^\infty(M).
$$

\begin{theorem}
\label{theorem-de-rham-characterization}
There is a unique sequence of operators
$$
C^\infty(M)\xrightarrow{d}\Lambda^1M
\xrightarrow{d}\Lambda^2M
\xrightarrow{d}\cdots
$$
such that $d$ is the ordinary differential
on functions, $d^2=0$, and, for a form
$\eta$ of degree $p$,
$$
d(\eta\wedge\xi)
=d\eta\wedge\xi
+(-1)^p\eta\wedge d\xi.
$$
\end{theorem}

\begin{definition}
\label{definition-de-rham-differential}
The operator $d$ in Theorem
\ref{theorem-de-rham-characterization} is
the {\it de Rham differential}.
\end{definition}

\begin{proposition}[Cartan formula]
\label{proposition-cartan-one-form}
For a one-form $\eta$ and vector fields
$X,Y$, one has
$$
\begin{aligned}
d\eta(X,Y)
&=\mathcal{L}_X(\eta(Y))\\
&\quad-\mathcal{L}_Y(\eta(X))\\
&\quad-\eta([X,Y]).
\end{aligned}
$$
\end{proposition}

\begin{proof}
The right hand side is alternating and
satisfies the Leibniz rule in $\eta$. It
vanishes when $\eta=df$, because
$$
\mathcal{L}_{[X,Y]}f
=\mathcal{L}_X\mathcal{L}_Yf
-\mathcal{L}_Y\mathcal{L}_Xf.
$$
These properties identify it with the de
Rham differential on one-forms.
\end{proof}

\medskip\noindent
{\bf Torsion.}

\begin{definition}
\label{definition-torsion-vector}
Let $\nabla$ be a connection on $TM$. Its
{\it torsion} is the tensor
$$
T_\nabla(X,Y)
=\nabla_XY-\nabla_YX-[X,Y].
$$
\end{definition}

The induced connection on $T^*M$ gives a
bundle map
$$
\operatorname{Alt}\circ\nabla:
\Lambda^1M\longrightarrow\Lambda^2M.
$$
Cartan's formula gives
$$
\begin{aligned}
(d\eta-&\operatorname{Alt}
\circ\nabla\eta)(X,Y)\\
&=\eta(T_\nabla(X,Y)).
\end{aligned}
$$
Thus the operator
$$
d-\operatorname{Alt}\circ\nabla
$$
on one-forms is dual to the usual
vector-valued torsion tensor. In particular,
it is $C^\infty(M)$-linear.

\begin{definition}
\label{definition-levi-civita}
Let $(M,g)$ be a Riemannian manifold. A
connection on $TM$ is
{\it Levi--Civita} if it is orthogonal and
torsion-free.
\end{definition}

\begin{theorem}
\label{theorem-fundamental-riemannian}
Every Riemannian manifold has a unique
Levi--Civita connection.
\end{theorem}

The slides recalled Gregorio
Ricci-Curbastro (1853--1925) and his former
student Tullio Levi-Civita (1873--1941).
Their joint work pioneered tensor calculus.
Ricci-Curbastro signed that publication with
the shortened surname Ricci, a convention
which has caused some later confusion.

\medskip\noindent
{\bf Linearized torsion.}

The space of all connections is affine, and
torsion is an affine map from this space to
$$
\Lambda^2M\otimes TM.
$$
If a connection is changed by
$$
A\in\Lambda^1M\otimes
\operatorname{End}(TM),
$$
its torsion changes by the
antisymmetrization
$$
\begin{aligned}
\operatorname{Alt}_{12}(A)(X,Y)
&=A(X)Y-A(Y)X.
\end{aligned}
$$

\begin{definition}
\label{definition-linearized-torsion}
The {\it linearized torsion} is the map
$$
\begin{aligned}
T_{\rm lin}:\Lambda^1M\otimes
\operatorname{End}(TM)
&\longrightarrow
\Lambda^2M\otimes TM,\\
A&\longmapsto
\operatorname{Alt}_{12}(A).
\end{aligned}
$$
\end{definition}

\begin{proposition}
\label{proposition-orthogonal-exists}
Every vector bundle with a scalar product
admits an orthogonal connection.
\end{proposition}

\begin{proof}
Choose a trivializing cover with local
orthonormal frames, and on each member choose
a connection preserving its frame. A
partition of unity gives an affine
combination of these local connections. The
result is a global orthogonal connection.
\end{proof}

\begin{proof}[Proof of Theorem
\ref{theorem-fundamental-riemannian}]
Choose an orthogonal connection $\nabla_0$.
By Proposition
\ref{proposition-affine-orthogonal}, the
space of orthogonal connections is affine
over
$$
\Lambda^1M\otimes\mathfrak{so}(TM).
$$
The metric identifies this bundle with
$$
\Lambda^1M\otimes\Lambda^2M.
$$
On this affine subspace, the linearized
torsion is
$$
\begin{aligned}
T_{\rm lin}:\Lambda^1M\otimes\Lambda^2M
&\longrightarrow
\Lambda^2M\otimes\Lambda^1M,\\
A&\longmapsto
\operatorname{Alt}_{12}(A).
\end{aligned}
$$
It is an isomorphism. The two bundles have
the same rank, so it is enough to prove
injectivity. In covariant form, an element
of the kernel is symmetric in its first two
arguments and skew-symmetric in its last
two. The proof is a permutation argument.
Let $\sigma$ be the cyclic permutation
$$
\sigma(x,y,z)=(y,z,x).
$$
The two symmetries give
$$
\sigma(\eta)=-\eta.
$$
Since $\sigma^3=\operatorname{Id}$, applying
the permutation three times gives
$$
\eta=-\eta,
$$
so $\eta=0$.

It follows that an orthogonal connection is
uniquely determined by its torsion. In
particular,
$$
\nabla
=\nabla_0
-T_{\rm lin}^{-1}(T_{\nabla_0})
$$
is the unique orthogonal connection with
zero torsion.
\end{proof}

\section{Lecture 9: the Bismut
connection}
\label{section-lecture-nine-bismut}

\noindent
This lecture was given on April 9, 2025.

\medskip\noindent
{\bf Intrinsic torsion.}

Let $\Phi$ be a tensor on a manifold and
let $\nabla$ preserve $\Phi$. Denote by
$$
\mathfrak{a}(M)
\subset\operatorname{End}(TM)
$$
the bundle of Lie algebras consisting of
the endomorphisms $A$ satisfying
$$
A(\Phi)=0.
$$
A second connection $\nabla_1$ preserves
$\Phi$ if and only if
$$
\nabla-\nabla_1
\in\Lambda^1M\otimes\mathfrak{a}(M).
$$
Thus the connections preserving $\Phi$
form an affine space over
$$
\Lambda^1M\otimes\mathfrak{a}(M).
$$

\begin{definition}
\label{definition-intrinsic-torsion-space}
The {\it space of intrinsic torsion} for
$\mathfrak{a}(M)$-valued connections is
the quotient bundle
$$
\mathcal{T}_{\mathfrak{a}}
=\frac{\Lambda^2M\otimes TM}
{\operatorname{Alt}_{12}
(\Lambda^1M\otimes\mathfrak{a}(M))}.
$$
\end{definition}

\begin{definition}
\label{definition-intrinsic-torsion}
Assume that $\mathfrak{a}(M)$ is a vector
bundle. The {\it intrinsic torsion} of a
connection preserving $\Phi$ is the image
of its torsion in
$\mathcal{T}_{\mathfrak{a}}$.
\end{definition}

\begin{theorem}
\label{theorem-intrinsic-torsion}
The intrinsic torsion is independent of
the choice of connection preserving
$\Phi$. There is a torsion-free connection
preserving $\Phi$ if and only if the
intrinsic torsion vanishes.
\end{theorem}

\begin{proof}
If two such connections differ by
$$
A\in\Lambda^1M\otimes\mathfrak{a}(M),
$$
their torsions differ by
$\operatorname{Alt}_{12}(A)$. This proves
the first assertion. The second follows by
changing a connection by a preimage of its
torsion under the linearized torsion map.
\end{proof}

\medskip\noindent
{\bf Decomposition of third tensor powers.}

For $1\leq i<j\leq n$, let
$$
\operatorname{Alt}_{ij},
\operatorname{Sym}_{ij}:
V^{\otimes n}\longrightarrow V^{\otimes n}
$$
be antisymmetrization and symmetrization in
the $i$-th and $j$-th factors.

\begin{proposition}
\label{proposition-third-tensor}
For every vector space $V$, one has
$$
\begin{aligned}
V^{\otimes3}
={}&\operatorname{Sym}^3V
\oplus\Lambda^3V\\
&\oplus\operatorname{Alt}_{12}
(V\otimes\operatorname{Sym}^2V)\\
&\oplus\operatorname{Sym}_{12}
(V\otimes\Lambda^2V).
\end{aligned}
$$
\end{proposition}

\begin{proof}
Every tensor $x$ satisfies
$$
x=\operatorname{Sym}_{ij}(x)
+\operatorname{Alt}_{ij}(x).
$$
First apply this to the last two factors to
obtain
$$
V^{\otimes3}
=V\otimes\operatorname{Sym}^2V
\oplus V\otimes\Lambda^2V.
$$
Then apply it to the first two factors.
The space
$$
\operatorname{Sym}_{12}
(V\otimes\operatorname{Sym}^2V)
$$
is symmetric in all three factors, while
$$
\operatorname{Alt}_{12}
(V\otimes\Lambda^2V)
$$
is antisymmetric in all three factors.
These are $\operatorname{Sym}^3V$ and
$\Lambda^3V$, respectively.
\end{proof}

\begin{lemma}[Corollary]
\label{lemma-symplectic-torsion-quotient}
There is a natural isomorphism
$$
\frac{\Lambda^2V\otimes V}
{\operatorname{Alt}_{12}
(V\otimes\operatorname{Sym}^2V)}
\cong\Lambda^3V.
$$
\end{lemma}

\begin{proof}
Antisymmetrizing the first two factors in
Proposition
\ref{proposition-third-tensor}
gives
$$
\begin{aligned}
\Lambda^2V\otimes V
={}&\Lambda^3V\\
&\oplus\operatorname{Alt}_{12}
(V\otimes\operatorname{Sym}^2V).
\end{aligned}
$$
\end{proof}

\medskip\noindent
{\bf Symplectic connections.}

For a connection on $T^*M$, define
$$
\tau_\nabla(\eta)
=\operatorname{Alt}(\nabla\eta)-d\eta.
$$
On one-forms this is the tensor dual to the
torsion. It extends to differential forms
as a graded derivation. Thus, if $\lambda$
has degree $p$, then
$$
\begin{aligned}
\tau_\nabla(\lambda\wedge\mu)
={}&\tau_\nabla(\lambda)\wedge\mu\\
&+(-1)^p\lambda\wedge
\tau_\nabla(\mu).
\end{aligned}
$$

\begin{definition}
\label{definition-almost-symplectic}
An {\it almost symplectic structure} is a
nondegenerate two-form.
\end{definition}

\begin{definition}
\label{definition-symplectic-connection}
Let $(M,\omega)$ be almost symplectic. A
{\it symplectic connection} is a connection
on $TM$ satisfying
$$
\nabla\omega=0.
$$
\end{definition}

Such connections exist: local connections
preserving symplectic frames can be joined
by an affine partition-of-unity argument.

\begin{lemma}
\label{lemma-symplectic-torsion-form}
Let $\nabla$ be a symplectic connection and
use $\omega$ to identify its torsion with
$$
T_\omega
\in\Lambda^2M\otimes\Lambda^1M.
$$
Then
$$
\operatorname{Alt}(T_\omega)
=-\frac{1}{2}d\omega.
$$
\end{lemma}

\begin{proof}
Since $\nabla\omega=0$, one has
$$
\tau_\nabla(\omega)=-d\omega.
$$
Apply the graded derivation identity to the
two covariant factors of $\omega$. The two
contractions with the torsion have the same
antisymmetrization. Each is
$\operatorname{Alt}(T_\omega)$, and hence
$$
\tau_\nabla(\omega)
=2\operatorname{Alt}(T_\omega).
$$
\end{proof}

\begin{theorem}
\label{theorem-almost-symplectic-torsion}
Let $(M,\omega)$ be almost symplectic. Use
$\omega$ to identify the torsion of a
symplectic connection with
$$
T_\omega
\in\Lambda^2M\otimes\Lambda^1M.
$$
Then
$$
\operatorname{Alt}_{123}(T_\omega)
=-\frac{1}{2}d\omega.
$$
Conversely, every tensor
$$
u\in\Lambda^2M\otimes\Lambda^1M
$$
satisfying
$$
\operatorname{Alt}_{123}(u)
=-\frac{1}{2}d\omega
$$
is the torsion of a symplectic connection.
\end{theorem}

\begin{proof}
The symplectic Lie algebra is identified by
$\omega$ with symmetric two-tensors:
$$
\mathfrak{sp}(TM)
\cong\operatorname{Sym}^2(\Lambda^1M).
$$
Under this identification, linearized
torsion becomes
$$
\begin{aligned}
\operatorname{Alt}_{12}:
\Lambda^1M\otimes
\operatorname{Sym}^2(\Lambda^1M)
\longrightarrow{}&
\Lambda^2M\otimes\Lambda^1M.
\end{aligned}
$$
Its kernel is
$\operatorname{Sym}^3(\Lambda^1M)$.
Lemma
\ref{lemma-symplectic-torsion-quotient}
therefore gives the exact sequence
$$
\begin{aligned}
0\longrightarrow{}&
\operatorname{Sym}^3(\Lambda^1M)\\
\longrightarrow{}&
\Lambda^1M\otimes
\operatorname{Sym}^2(\Lambda^1M)\\
\xrightarrow{\operatorname{Alt}_{12}}{}&
\Lambda^2M\otimes\Lambda^1M\\
\xrightarrow{\operatorname{Alt}_{123}}{}&
\Lambda^3M\longrightarrow0.
\end{aligned}
$$
Thus $\Lambda^3M$ is the space of
intrinsic torsion of symplectic
connections. Lemma
\ref{lemma-symplectic-torsion-form}
identifies the intrinsic torsion with
$-d\omega/2$, proving both assertions.
\end{proof}

\medskip\noindent
{\bf Torsion of unitary connections.}

\begin{proposition}
\label{proposition-unitary-torsion-types}
Let $(M,I,\omega)$ be a Hermitian complex
manifold and let $\nabla$ preserve $I$ and
$\omega$. Use the Riemannian metric to
regard the torsion as an element of
$$
\Lambda^2M\otimes\Lambda^1M.
$$
Then it belongs to
$$
\begin{aligned}
\mathcal{W}={}&
\Lambda^{2,0}(M)\otimes\Lambda^{0,1}(M)\\
&\oplus\Lambda^{0,2}(M)
\otimes\Lambda^{1,0}(M)\\
&\oplus\Lambda^{1,1}(M)
\otimes\Lambda^1M.
\end{aligned}
$$
\end{proposition}

\begin{proof}
Integrability of $I$ implies
$$
[T^{1,0}M,T^{1,0}M]\subset T^{1,0}M.
$$
Since $\nabla I=0$, one also has
$$
\nabla_X(T^{1,0}M)\subset T^{1,0}M
$$
for every vector field $X$. It follows that
the torsion lies in
$$
\begin{aligned}
&\Lambda^{2,0}(M)\otimes T^{1,0}M\\
&\quad\oplus\Lambda^{0,2}(M)
\otimes T^{0,1}M\\
&\quad\oplus\Lambda^{1,1}(M)\otimes TM.
\end{aligned}
$$
The Riemannian metric has type $(1,1)$, so
it identifies $T^{1,0}M$ with
$\Lambda^{0,1}(M)$ and $T^{0,1}M$ with
$\Lambda^{1,0}(M)$.
\end{proof}

The slides recalled Jean-Michel Bismut,
who was born in 1948.

\begin{theorem}[Bismut]
\label{theorem-bismut-connection}
Let $(M,I,\omega)$ be a Hermitian complex
manifold. There is a unique connection
$\nabla$ preserving $I$ and $\omega$ whose
torsion is a three-form. Its torsion is
$$
T_\nabla=-\frac{1}{2}I(d\omega).
$$
\end{theorem}

\begin{definition}
\label{definition-bismut-connection}
The connection in Theorem
\ref{theorem-bismut-connection} is the
{\it Bismut connection}.
\end{definition}

\begin{proof}
There are two identifications
$$
TM\longrightarrow\Lambda^1M,
$$
one given by $g$ and the other by $\omega$.
Write $T_g$ and $T_\omega$ for the two
covariant forms of the same torsion. Define
$I_3$ by
$$
I_3(x\otimes y\otimes z)
=x\otimes y\otimes I(z).
$$
Then
$$
I_3(T_g)=T_\omega.
$$
By Theorem
\ref{theorem-almost-symplectic-torsion},
$$
\operatorname{Alt}(I_3(T_g))
=-\frac{1}{2}d\omega.
$$
The left hand side is independent of the
unitary connection. Its linearization
therefore vanishes:
$$
\operatorname{Alt}\circ I_3\circ
T_{\rm lin}=0
$$
on
$$
\Lambda^1M\otimes\mathfrak{u}(TM).
$$

By Proposition
\ref{proposition-unitary-torsion-types},
unitary torsion belongs to $\mathcal{W}$.
The linearized torsion map
$$
T_{\rm lin}:
\Lambda^1M\otimes\mathfrak{u}(TM)
\longrightarrow\mathcal{W}
$$
is injective by the same permutation
argument used for the Levi--Civita
connection. Dimension counting gives the
exact sequence
$$
\begin{aligned}
0\longrightarrow{}&
\Lambda^1M\otimes\mathfrak{u}(TM)\\
\xrightarrow{T_{\rm lin}}{}&
\mathcal{W}\\
\xrightarrow{\operatorname{Alt}\circ I_3}{}&
\Lambda^{2,1}(M)\oplus\Lambda^{1,2}(M)\\
\longrightarrow{}&0.
\end{aligned}
$$
Surjectivity of the second map follows by
antisymmetrizing tensors in
$I_3(\mathcal{W})$.

Let
$$
\mathcal{U}
=\Lambda^{2,1}(M)\oplus\Lambda^{1,2}(M)
$$
be the three-forms contained in
$\mathcal{W}$. On three-forms one has
$$
\operatorname{Alt}(I_3(\eta))
=\mathcal{W}_I(\eta),
$$
where the Weil operator is given by
$$
\begin{aligned}
\mathcal{W}_I(\eta)(x,y,z)
={}&\eta(Ix,y,z)\\
&+\eta(x,Iy,z)\\
&+\eta(x,y,Iz).
\end{aligned}
$$
Its restriction from $\mathcal{U}$ to
itself is bijective. Consequently there is
a unique $\sigma\in\mathcal{U}$ satisfying
$$
\operatorname{Alt}(I_3(\sigma))
=-\frac{1}{2}d\omega.
$$

Choose a unitary connection $\nabla_0$ and
let $T_0\in\mathcal{W}$ be its torsion.
Exactness gives a unique
$$
A\in\Lambda^1M\otimes\mathfrak{u}(TM)
$$
such that
$$
T_{\rm lin}(A)+T_0=\sigma.
$$
The connection $\nabla_0+A$ is unitary and
has totally antisymmetric torsion. The Weil
operator calculation on forms of types
$(2,1)$ and $(1,2)$ gives
$$
\sigma=-\frac{1}{2}I(d\omega).
$$
\end{proof}

\begin{lemma}[Corollary]
\label{lemma-kahler-levi-civita-unitary}
On a K"ahler manifold, the Levi--Civita
connection satisfies
$$
\nabla I=0,
\qquad
\nabla\omega=0.
$$
\end{lemma}

\begin{proof}
If $d\omega=0$, the Bismut connection has
zero torsion. It is orthogonal, so uniqueness
of the Levi--Civita connection identifies
the two connections.
\end{proof}

There is also a direct intrinsic-torsion
proof. For a unitary connection $\nabla$ on
a K"ahler manifold, Theorem
\ref{theorem-almost-symplectic-torsion}
gives
$$
\operatorname{Alt}(I_3(T_\nabla))=0.
$$
The exact sequence in the proof of Theorem
\ref{theorem-bismut-connection} implies
that $T_\nabla$ is in the image of
$T_{\rm lin}$. Hence
$$
\nabla-T_{\rm lin}^{-1}(T_\nabla)
$$
is a torsion-free unitary connection and
therefore is the Levi--Civita connection.

\section{Lecture 9-bis: harmonic forms and
cohomology}
\label{section-lecture-nine-bis-harmonic}

\noindent
In this lecture I learnt that
we can define harmonic forms on any
chain of vector spaces equipped with
an inner product. Define $d^*$
as always, satisfying 
$\langle d \alpha, \beta\rangle
=
\langle \alpha, d^*\beta\rangle$,
and the Laplacian, as always by
$d d^* + d^* d$.
The key observation is that
$$
H^k = \frac{\ker d}{\operatorname{Im} d}
\simeq
\ker d \cap \operatorname{Im} (d)^\perp
\simeq
\ker d \cap \ker d^*
\simeq
\ker \Delta,
$$
which says that evert cohomology
class is represented by a harmonic
vector.

\medskip\noindent
{\bf Laplacians on a complex.}

\begin{definition}
\label{definition-complex-vector-spaces}
A {\it complex of vector spaces} is a
sequence of vector spaces and linear maps
$$
\cdots\xrightarrow{d}C^{i-1}
\xrightarrow{d}C^i
\xrightarrow{d}C^{i+1}
\xrightarrow{d}\cdots
$$
satisfying $d^2=0$. It is {\it finite} if
only finitely many of its vector spaces are
nonzero.
\end{definition}

\begin{definition}
\label{definition-cohomology-complex}
The {\it cohomology} of $(C^*,d)$ in degree
$i$ is
$$
H^i(C^*)
=\frac{\ker(d|_{C^i})}{d(C^{i-1})}.
$$
\end{definition}

\begin{definition}
\label{definition-laplacian-complex}
Suppose that every $C^i$ has a positive
definite inner product. If $d^*$ denotes
the adjoint of $d$, the {\it Laplacian} is
$$
\Delta_d=dd^*+d^*d.
$$
\end{definition}

The Laplacian commutes with $d$, and hence
is an endomorphism of the complex.

\begin{definition}
\label{definition-harmonic-vector}
A vector in $C^i$ is {\it harmonic} if it
belongs to $\ker\Delta_d$.
\end{definition}

For every $x\in C^i$, one has
$$
\begin{aligned}
\left\langle{}x,\Delta_dx
\right\rangle{}
={}&\left\langle{}dx,dx
\right\rangle{}\\
&+\left\langle{}d^*x,d^*x
\right\rangle{}.
\end{aligned}
$$
Therefore
$$
\ker\Delta_d=\ker d\cap\ker d^*.
$$

The adjoint identities give
$$
\ker d=(\operatorname{im}d^*)^\perp,
\qquad
\ker d^*=(\operatorname{im}d)^\perp.
$$
When the relevant images are closed, they
also give
$$
\operatorname{im}d
=(\ker d^*)^\perp,
\qquad
\operatorname{im}d^*
=(\ker d)^\perp.
$$

\begin{lemma}[Corollary]
\label{lemma-harmonic-representatives}
Suppose that the vector spaces $C^i$ are
finite-dimensional. More generally, suppose
that the images of the differentials are
closed in the relevant Hilbert spaces. Then
every cohomology class has a unique harmonic
representative, and
$$
H^*(C^*)\cong\ker\Delta_d.
$$
\end{lemma}

\begin{proof}
Orthogonal decomposition gives
$$
\begin{aligned}
\ker d
&=\operatorname{im}d
&\quad\oplus
(\ker d\cap(\operatorname{im}d)^\perp)\\
&=\operatorname{im}d
\oplus(\ker d\cap\ker d^*)\\
&=\operatorname{im}d
\oplus\ker\Delta_d.
\end{aligned}
$$
The quotient of $\ker d$ by
$\operatorname{im}d$ is therefore naturally
isomorphic to $\ker\Delta_d$.
\end{proof}

\medskip\noindent
{\bf Laplacians on differential forms.}

A metric $g$ on a vector space induces a
metric on every tensor power by
$$
\begin{aligned}
&g(x_1\otimes\cdots\otimes x_k,
y_1\otimes\cdots\otimes y_k)\\
&\qquad
=\prod_{j=1}^kg(x_j,y_j).
\end{aligned}
$$
Consequently a Riemannian metric on $M$
induces the $L^2$ inner product
$$
\left\langle{}\alpha,\beta
\right\rangle{}_{L^2}
=\int_M g(\alpha,\beta)
\operatorname{Vol}_M
$$
on differential forms.

\begin{definition}
\label{definition-laplacian-forms}
The {\it Laplacian on differential forms}
on a Riemannian manifold is
$$
\Delta=dd^*+d^*d.
$$
\end{definition}

\begin{theorem}[Hodge theorem]
\label{theorem-hodge-spectral}
Let $M$ be a compact Riemannian manifold.
The Hilbert space $L^2(\Lambda^*M)$ has an
orthonormal basis of eigenforms of
$\Delta$. Every eigenspace is
finite-dimensional, and the spectrum is
discrete. The inverse of $\Delta$ on
$\operatorname{im}\Delta$ is continuous in
the $L^2$ topology.
\end{theorem}

\begin{theorem}[Elliptic regularity]
\label{theorem-laplacian-regularity}
Every $L^2$ eigenform of the Laplacian on a
compact Riemannian manifold is smooth.
\end{theorem}

The analogous spectral and regularity
statements hold for the Dolbeault Laplacian
on a compact Hermitian manifold.

The slides recalled Fritz Noether
(1884--1941), the brother of Emmy Noether.

\medskip\noindent
{\bf De Rham cohomology.}

\begin{definition}
\label{definition-de-rham-cohomology}
The {\it de Rham cohomology} of $M$ is
$$
H^i(M)
=\frac{\ker(d|_{\Lambda^iM})}
{d(\Lambda^{i-1}M)}.
$$
\end{definition}

\begin{definition}
\label{definition-harmonic-form}
A differential form $\alpha$ is
{\it harmonic} if
$$
\Delta\alpha=0.
$$
The space of harmonic $i$-forms is denoted
by $\mathcal{H}^i(M)$.
\end{definition}

A harmonic form is both $d$-closed and
$d^*$-closed. The natural map
$$
\mathcal{H}^i(M)\longrightarrow H^i(M)
$$
is injective. Indeed, if
$\alpha=d\beta$ is harmonic, then
$$
\begin{aligned}
\left\langle{}\alpha,\alpha
\right\rangle{}
&=\left\langle{}\alpha,d\beta
\right\rangle{}\\
&=\left\langle{}d^*\alpha,\beta
\right\rangle{}\\
&=0.
\end{aligned}
$$

\begin{theorem}[Hodge decomposition]
\label{theorem-hodge-de-rham}
For a compact Riemannian manifold, the
natural map
$$
\mathcal{H}^i(M)\longrightarrow H^i(M)
$$
is an isomorphism.
\end{theorem}

\begin{proof}
The identities $d^2=0$ and $(d^*)^2=0$
imply
$$
[d,\Delta]=0.
$$
Thus $d$ preserves every eigenspace of the
Laplacian. Let
$$
\Lambda^*(M)
=\mathop{\widehat{\bigoplus}}_\lambda
\Lambda^*_{\lambda}(M)
$$
be the eigenspace decomposition. Each
$\Lambda^*_{\lambda}(M)$ is a complex under
$d$.

Suppose that $\lambda\neq0$ and that
$\eta\in\Lambda^*_{\lambda}(M)$ is closed.
Then
$$
\begin{aligned}
\lambda\eta
&=dd^*\eta+d^*d\eta\\
&=dd^*\eta.
\end{aligned}
$$
It follows that
$$
\eta=d(\lambda^{-1}d^*\eta).
$$
Hence the nonzero eigenspaces have no
cohomology. The zero eigenspace is
$\mathcal{H}^*(M)$, which proves the
claim.
\end{proof}

Poincar\'e duality follows by applying the
Hodge star to harmonic forms.

\medskip\noindent
{\bf Dolbeault cohomology.}

\begin{definition}
\label{definition-dolbeault-cohomology}
The {\it Dolbeault cohomology} of a complex
manifold $M$ is
$$
H^{p,q}_{\bar\partial}(M)
=\frac{\ker(\bar\partial|
_{\Lambda^{p,q}(M)})}
{\bar\partial\Lambda^{p,q-1}(M)}.
$$
\end{definition}

On a Hermitian manifold, define
$$
\Delta_{\bar\partial}
=\bar\partial\bar\partial^*
+\bar\partial^*\bar\partial
$$
and
$$
\mathcal{H}^{p,q}_{\bar\partial}(M)
=\ker(\Delta_{\bar\partial}|
_{\Lambda^{p,q}(M)}).
$$

\begin{theorem}
\label{theorem-hodge-dolbeault}
Let $M$ be a compact Hermitian complex
manifold. The natural map
$$
\mathcal{H}^{p,q}_{\bar\partial}(M)
\longrightarrow
H^{p,q}_{\bar\partial}(M)
$$
is an isomorphism.
\end{theorem}

\begin{proof}
The identities
$$
\bar\partial^2=0,
\qquad
(\bar\partial^*)^2=0
$$
give
$$
[\bar\partial,
\Delta_{\bar\partial}]=0.
$$
Therefore $\bar\partial$ preserves the
eigenspaces of $\Delta_{\bar\partial}$.

Let $\lambda\neq0$, and let $\eta$ be a
$\bar\partial$-closed eigenform with
eigenvalue $\lambda$. Then
$$
\begin{aligned}
\lambda\eta
&=\bar\partial\bar\partial^*\eta
+\bar\partial^*\bar\partial\eta\\
&=\bar\partial\bar\partial^*\eta.
\end{aligned}
$$
Consequently
$$
\eta
=\bar\partial
(\lambda^{-1}\bar\partial^*\eta).
$$
The nonzero eigenspaces contribute no
Dolbeault cohomology. The zero eigenspace is
$\mathcal{H}^{p,q}_{\bar\partial}(M)$.
\end{proof}

\end{document}
